\chapter{O Sacramento da Eucaristia}

A sagrada Eucaristia completa a iniciação cristã. Aqueles que foram elevados à dignidade do sacerdócio real pelo Batismo e configurados mais profundamente com Cristo pela Confirmação participam, por meio da Eucaristia, com toda a comunidade, no próprio sacrifício do Senhor. O Catecismo da Igreja Católica trata deste sacramento nos parágrafos 1322 a 1419, expondo-o como o coração e o cume de toda a vida da Igreja. Não é um sacramento entre outros, mas o sacramento que resume e sustenta todos; nele a Igreja não apenas recebe o dom de Cristo, mas ela própria é constituída como evento sacramental, oferecida ao Pai como sacrifício de louvor e de ação de graças.

\section{Fonte e cume da vida eclesial (1322--1327)}

A sagrada Eucaristia completa a iniciação cristã e ocupa um lugar único entre os sacramentos: é ``fonte e cume de toda a vida cristã''. A razão desta primazia está em que ``na santíssima Eucaristia está contido todo o tesouro espiritual da Igreja, isto é, o próprio Cristo, nossa Páscoa''. Os demais sacramentos, assim como todos os ministérios eclesiais e obras de apostolado, estão vinculados à Eucaristia e a ela se ordenam. Na Eucaristia, a Igreja recebe aquilo pelo qual ela existe: o próprio Cristo, e nada além d'Ele.

A comunhão de vida com Deus e a unidade do povo de Deus --- pelas quais a Igreja é o que é --- são significadas e realizadas pela Eucaristia. Nela se encontra o cume da ação pela qual Deus, em Cristo, santifica o mundo, e ao mesmo tempo o cume do culto que, no Espírito Santo, os homens prestam a Cristo e, por Ele, ao Pai. Celebrando-a, a Igreja une-se desde já à liturgia do céu e antecipa a vida eterna, quando ``Deus for tudo em todos''. Por isso, a Eucaristia não é apenas uma prática de devoção entre muitas, mas a ação mais alta que a Igreja pode realizar nesta vida.

O Senhor instituiu a Eucaristia ``na última ceia, na noite em que foi entregue'', para perpetuar pelo decurso dos séculos, até voltar, o sacrifício da Cruz, confiando à Igreja, sua Esposa amada, o memorial da sua morte e ressurreição: ``sacramento de piedade, sinal de unidade, vínculo de caridade, banquete pascal em que se recebe Cristo, a alma se enche de graça e nos é dado o penhor da glória futura''. Esta instituição não foi um gesto isolado, mas o ato pelo qual Cristo entrou para sempre na história da Igreja como o Senhor que Se oferece, que alimenta, que une e que envia.

A centralidade da Eucaristia na vida cristã implica consequências concretas para toda a catequese, para a espiritualidade e para a missão. Não é possível compreender a Igreja como comunidade, o sacerdócio comum dos fiéis, a vocação à santidade ou o compromisso apostólico prescindindo da Eucaristia. Toda a existência cristã gravita em torno deste centro: é alimentada por ele, purificada nele, enviada a partir dele. A catequese eucarística consiste, por isso, não em explicar um rito, mas em introduzir os fiéis no interior do mistério pelo qual Cristo os incorpora ao seu próprio oferecer-se ao Pai.

\section{Os nomes da Eucaristia (1328--1332)}

A riqueza inesgotável deste sacramento exprime-se nos diferentes nomes que lhe são dados. Chama-se ``Eucaristia'' porque é ação de graças a Deus; as palavras gregas \emph{eucharistein} e \emph{eulogein} evocam as bênçãos judaicas que proclamavam as obras de Deus --- criação, redenção e santificação. É chamada ``Ceia do Senhor'' porque se trata da ceia que o Senhor comeu com os discípulos na véspera da sua Paixão e da antecipação do banquete nupcial do Cordeiro na Jerusalém celeste. É designada ``Fração do Pão'' porque este gesto, próprio das refeições judaicas, foi utilizado por Jesus quando abençoava e distribuía o pão, sobretudo na última Ceia; por este gesto os discípulos O reconheceram depois da ressurreição, e com esta expressão os primeiros cristãos designavam as suas assembleias eucarísticas.

Chama-se ``Assembleia Eucarística'' porque é celebrada em assembleia dos fiéis, expressão visível da Igreja. É chamada ``Memorial da Paixão e Ressurreição do Senhor''. É o ``Santo Sacrifício'' que atualiza o único sacrifício de Cristo Salvador e inclui a oferenda da Igreja; nesse sentido, diz-se também ``sacrifício de louvor'', ``sacrifício espiritual'', ``sacrifício puro e santo'', que completa e ultrapassa todos os sacrifícios da Antiga Aliança. É chamada ``Santa e Divina Liturgia'', porque toda a liturgia da Igreja encontra o seu centro e a expressão mais densa neste sacramento. Fala-se igualmente do ``Santíssimo Sacramento'', porque é o sacramento dos sacramentos, e sob este nome se designam também as espécies eucarísticas guardadas no sacrário.

Chama-se ``Comunhão'' porque, por este sacramento, nos unimos a Cristo, que nos torna participantes do seu Corpo e do seu Sangue para formarmos um só Corpo. Designa-se ainda ``as coisas santas'' --- é este o sentido primário da ``comunhão dos santos'' de que fala o Símbolo dos Apóstolos ---, ``pão dos anjos'', ``pão do céu'', ``remédio da imortalidade'', ``viático''. Finalmente, chama-se ``Santa Missa'' porque a liturgia em que se realiza o mistério da salvação termina com o envio dos fiéis --- \emph{missio} --- para que vão cumprir a vontade de Deus na sua vida quotidiana. A Missa é, pois, ao mesmo tempo, o ponto de chegada de toda a vida cristã e o ponto de partida da sua missão no mundo.

Esta pluralidade de nomes não é redundância, mas riqueza teológica: cada título ilumina uma dimensão do mesmo mistério inexaurível. Nenhum nome esgota o que a Eucaristia é; juntos, eles circunscrevem um horizonte de sentido que só pode ser penetrado pela fé viva e pela participação ativa na própria celebração. A catequese eucarística deve, por isso, introduzir os fiéis nesta linguagem simbólica, não para que memorizem definições, mas para que aprendam a habitar o mistério com inteligência e amor.

\section{A Eucaristia na economia da salvação (1333--1344)}

No centro da celebração eucarística estão o pão e o vinho que, pelas palavras de Cristo e pela invocação do Espírito Santo, se tornam o Corpo e o Sangue do mesmo Cristo. Tornando-se misteriosamente o Corpo e o Sangue de Cristo, os sinais do pão e do vinho continuam a significar a bondade da criação; por isso, no ofertório, a Igreja dá graças ao Criador pelo pão e pelo vinho, fruto ``do trabalho do homem'' mas primeiramente ``fruto da terra'' e ``da videira''. Na Antiga Aliança, o pão e o vinho eram oferecidos em sacrifício entre as primícias da terra; os pães ázimos da Páscoa judaica comemoravam a partida libertadora do Egito; e o ``cálice de bênção'' ao fim da ceia pascal judaica acrescentava à alegria do vinho uma dimensão escatológica --- a expectativa messiânica do restabelecimento de Jerusalém. Jesus instituiu a sua Eucaristia dando um sentido novo e definitivo a esta bênção.

Os milagres da multiplicação dos pães, quando o Senhor disse a bênção, partiu e distribuiu os pães pelos discípulos para alimentar a multidão, prefiguram a superabundância deste pão único da Eucaristia. O sinal da água transformada em vinho em Caná já anuncia a ``Hora'' da glorificação de Jesus e manifesta o cumprimento do banquete das núpcias no Reino do Pai. O primeiro anúncio da Eucaristia, no discurso do Pão da Vida em Cafarnaum, dividiu os discípulos, tal como o anúncio da Paixão os escandalizou: ``Estas palavras são insuportáveis! Quem as pode escutar?'' A cruz e a Eucaristia são o mesmo mistério e não cessam de ser ocasião de divisão; mas só Cristo tem ``palavras de vida eterna'', e acolher na fé o dom da sua Eucaristia é acolhê-L'O a Ele próprio.

Tendo amado os seus, o Senhor amou-os até ao fim. Sabendo que era chegada a hora de partir deste mundo para regressar ao Pai, no decurso duma refeição lavou-lhes os pés e deu-lhes o mandamento do amor. Para lhes deixar uma garantia deste amor, para jamais se afastar dos seus e para os tornar participantes da sua Páscoa, instituiu a Eucaristia como memorial da sua morte e ressurreição, ordenando aos seus Apóstolos que a celebrassem até ao seu regresso, constituindo-os ``sacerdotes do Novo Testamento''. Jesus escolheu a Páscoa para cumprir o que tinha anunciado: na última Ceia, tomou o pão e, dando graças, partiu-o e deu-lho dizendo: ``Isto é o Meu corpo, que vai ser entregue por vós. Fazei isto em memória de Mim''; e da mesma forma o cálice: ``Este cálice é a Nova Aliança no meu sangue, que vai ser derramado por vós''.

Ao ordenar que repetissem os seus gestos e palavras ``até que Ele venha'', Jesus não pede somente que se lembrem d'Ele e do que fez. Tem em vista a celebração litúrgica, pelos Apóstolos e seus sucessores, do memorial de Cristo --- da sua vida, morte, ressurreição e intercessão junto do Pai. Desde o dia de Pentecostes, a Igreja foi fiel a esta ordem do Senhor: a comunidade de Jerusalém era assídua ``ao ensinamento dos Apóstolos, à união fraterna, à Fração do Pão e às orações''. Era sobretudo ``no primeiro dia da semana'' que os cristãos se reuniam ``para partir o Pão''. Desde então até hoje, a celebração da Eucaristia perpetuou-se de maneira que a encontramos em toda a parte na Igreja com a mesma estrutura fundamental. Ela continua a ser o centro da vida da Igreja, pois pela Eucaristia, de celebração em celebração, o Povo de Deus em peregrinação ``avança pela porta estreita da Cruz'' para o banquete celeste.

\section{A celebração litúrgica da Eucaristia (1345--1355)}

Desde o século II, temos o testemunho de São Justino Mártir sobre as grandes linhas do desenrolar da celebração eucarística, que permaneceram as mesmas até hoje em todas as grandes famílias litúrgicas. Por volta do ano 155, descrevendo a celebração ao imperador pagão, Justino indica: leem-se as memórias dos Apóstolos e os escritos dos Profetas; aquele que preside exorta à imitação dessas obras; levantam-se todos, fazem orações, dão o ósculo da paz; apresentam-se ao presidente pão e uma taça de água e vinho; ele faz subir louvor e glória ao Pai pelo nome do Filho e do Espírito Santo, e dá graças longamente; o povo responde ``Ámen''; e por fim os diáconos distribuem a todos os presentes o pão e o vinho ``eucaristizados'', levando-os também aos ausentes.

A liturgia eucarística processa-se segundo uma estrutura fundamental desdobrada em dois grandes momentos que formam uma unidade: a liturgia da Palavra, com as leituras, a homilia e a oração universal; e a liturgia eucarística, com a apresentação do pão e do vinho, a ação de graças consecratória e a Comunhão. Liturgia da Palavra e liturgia eucarística constituem juntas ``um só e mesmo ato de culto'', porque a mesa posta para nós na Eucaristia é, ao mesmo tempo, a da Palavra de Deus e a do Corpo do Senhor. É o mesmo dinamismo da refeição pascal de Jesus Ressuscitado com os discípulos de Emaús: enquanto caminhavam, Ele explicava-lhes as Escrituras; e depois, pondo-Se à mesa com eles, tomou o pão, proferiu a bênção, partiu-o e deu-lho.

A ação de graças consecratória --- a anáfora --- é o coração e o cume da celebração. No prefácio, a Igreja dá graças ao Pai, por Cristo, no Espírito Santo, por todas as suas obras: criação, redenção e santificação; toda a comunidade une as suas vozes ao louvor incessante da Igreja celeste, cantando o ``Santo, Santo, Santo''. Na \emph{epiclese}, a Igreja pede ao Pai que envie o seu Espírito Santo sobre o pão e o vinho para que se tornem o Corpo e o Sangue de Cristo, e para que os participantes sejam um só Corpo e um só Espírito. Segue-se a narração da instituição, em que a força das palavras e da ação de Cristo e o poder do Espírito Santo tornam sacramentalmente presentes, sob as espécies do pão e do vinho, o Corpo e o Sangue de Cristo, o seu sacrifício oferecido na Cruz de uma vez por todas. Na anamnese, a Igreja faz memória da Paixão, Ressurreição e retorno glorioso de Cristo, apresentando ao Pai a oferenda do seu Filho; nas intercessões, manifesta que a Eucaristia é celebrada em comunhão com toda a Igreja do céu e da terra, dos vivos e dos defuntos, e na comunhão com os pastores da Igreja.

A Comunhão, precedida da Oração do Senhor e da Fração do Pão, é o momento em que os fiéis recebem ``o pão do céu'' e ``o cálice da salvação'', o Corpo e o Sangue de Cristo que Se entregou ``para a vida do mundo''. São Justino testemunha que este alimento é chamado Eucaristia e que ninguém pode tomar parte nela se não acreditar na verdade do que se ensina, se não recebeu o batismo e se não vive segundo os preceitos de Cristo. A Fração do Pão evoca a última Ceia, mas remete também a todos os episódios em que Jesus partiu o pão --- na multiplicação dos pães, em Emaús, no Cenáculo. Neste gesto simples está inscrita toda a teologia da Eucaristia: o único pão partido é Cristo, e todos os que dele comem formam um só corpo n'Ele.

\section{O sacrifício sacramental: ação de graças, memorial, presença (1356--1381)}

Se os cristãos celebram a Eucaristia desde as origens e sob uma forma que na sua substância não mudou, é porque sabem que estão ligados pela ordem do Senhor: ``Fazei isto em memória de Mim''. Esta ordem cumprimo-la celebrando o memorial do seu sacrifício; e fazendo-o, oferecemos ao Pai o que Ele próprio nos deu --- os dons da sua criação, o pão e o vinho, transformados pelo poder do Espírito Santo e pelas palavras de Cristo no Corpo e no Sangue do mesmo Cristo. A Eucaristia é assim, inseparavelmente, ação de graças e louvor ao Pai, memorial sacrificial de Cristo e do seu Corpo, e presença de Cristo pelo poder da sua Palavra e do seu Espírito.

Como ação de graças, a Eucaristia é o sacrifício de louvor pelo qual a Igreja canta a glória de Deus em nome de toda a criação; este sacrifício só é possível por Cristo, que une os fiéis à sua pessoa, ao seu louvor e à sua intercessão. A Eucaristia é o memorial da Páscoa de Cristo --- a atualização e a oferenda sacramental do seu único sacrifício, na liturgia da Igreja que é o seu Corpo. No sentido que lhe dá a Sagrada Escritura, o memorial não é somente a lembrança de acontecimentos passados, mas a proclamação das maravilhas que Deus fez e a sua tornada presente: ``Todas as vezes que no altar se celebra o sacrifício da Cruz, no qual "Cristo, nossa Páscoa, foi imolado", realiza-se a obra da nossa redenção''. O sachrifício de Cristo e o sacrifício da Eucaristia são um único sacrifício: ``É uma só e mesma vítima, e Aquele que agora Se oferece pelo ministério dos sacerdotes é o mesmo que outrora Se ofereceu a Si mesmo na Cruz; só a maneira de oferecer é que é diferente''.

A Eucaristia é igualmente o sacrifício da Igreja. A Igreja, que é o Corpo de Cristo, participa na oblação da sua Cabeça: com Ele, ela própria é oferecida integralmente. Ela une-se à intercessão de Cristo junto do Pai em favor de todos os homens; a vida dos fiéis, o seu louvor, o seu sofrimento, a sua oração, o seu trabalho unem-se aos de Cristo e à sua oblação total, adquirindo assim um novo valor redentor. A Eucaristia é também oferecida pelos fiéis defuntos ``que morreram em Cristo e que não estão ainda de todo purificados'', para que possam entrar na luz e na paz de Cristo. Santo Agostinho resumiu admiravelmente: ``Toda esta cidade resgatada, oferecida a Deus como sacrifício universal pelo Sumo-Sacerdote que Se ofereceu a Si mesmo na Paixão para fazer de nós o Corpo de tal Cabeça --- este sacrifício a Igreja não cessa de o renovar no sacramento do altar''.

O modo da presença de Cristo sob as espécies eucarísticas é único e eleva a Eucaristia acima de todos os sacramentos. ``No santíssimo sacramento da Eucaristia estão contidos, verdadeira, real e substancialmente, o Corpo e o Sangue, conjuntamente com a alma e a divindade de nosso Senhor Jesus Cristo e, por conseguinte, Cristo completo.'' Esta presença chama-se ``real'' não a título exclusivo --- como se as outras presenças não fossem reais ---, mas por excelência, porque é substancial. A presença de Cristo começa no momento da consagração e dura enquanto as espécies eucarísticas subsistirem. O Concílio de Trento resume a fé católica declarando que, ``pela consagração do pão e do vinho, opera-se a conversão de toda a substância do pão na substância do Corpo de Cristo, e de toda a substância do vinho na substância do seu Sangue''; a esta mudança, a Igreja chama, de modo conveniente e apropriado, ``transubstanciação''. A Igreja Católica sempre prestou e continua a prestar culto de adoração à Eucaristia, não só durante a Missa, mas também fora da sua celebração, conservando as hóstias consagradas e levando-as em procissão.

\section{O banquete pascal: a Comunhão e seus frutos (1382--1401)}

A Missa é ao mesmo tempo e inseparavelmente o memorial sacrificial em que se perpetua o sacrifício da Cruz e o banquete sagrado da Comunhão do Corpo e Sangue do Senhor. O Senhor dirige-nos um convite insistente: ``Em verdade, em verdade vos digo: se não comerdes a Carne do Filho do Homem e não beberdes o seu Sangue, não tereis a vida em vós''. Para responder dignamente a este convite, é preciso preparação: São Paulo exorta a um exame de consciência, pois ``quem comer o pão ou beber do cálice do Senhor indignamente será réu do Corpo e do Sangue do Senhor''. Aquele que tiver consciência dum pecado grave deve receber o sacramento da Reconciliação antes de se aproximar da Comunhão. A atitude corporal --- gestos, traje --- deve traduzir o respeito, a solenidade e a alegria deste momento em que Cristo Se torna nosso hóspede.

Receber a Eucaristia na Comunhão traz como fruto principal a união íntima com Cristo Jesus: ``Quem come a minha Carne e bebe o meu Sangue permanece em Mim e Eu nele''. Como o alimento material restaura as forças perdidas, assim a Eucaristia fortalece a caridade que, na vida quotidiana, tende a enfraquecer-se. A Comunhão afasta-nos do pecado: apaga os pecados veniais e preserva dos pecados mortais futuros, pois quanto mais participarmos na vida de Cristo, mais difícil será romper com Ele pelo pecado mortal. Pela Eucaristia, os que a recebem ficam mais estreitamente unidos a Cristo, e Cristo une assim todos os fiéis num só Corpo: ``O cálice da bênção que abençoamos, não é comunhão com o Sangue de Cristo? O pão que partimos, não é comunhão com o Corpo de Cristo? Uma vez que há um único pão, nós, embora muitos, somos um só corpo, porque participamos desse único pão''.

A Eucaristia compromete-nos também com os pobres: para receber verdadeiramente o Corpo e o Sangue de Cristo entregue por nós, é preciso reconhecer Cristo nos mais pobres, seus irmãos: ``Saboreaste o Sangue do Senhor e não reconheces sequer o teu irmão. Desonras esta mesa, se não julgas digno de partilhar o teu alimento aquele que foi julgado digno de tomar parte nesta mesa''. A Eucaristia é também o centro da questão da unidade dos cristãos: ``Ó sacramento da piedade, ó sinal de unidade, ó vínculo de caridade!'', exclamava Santo Agostinho. Tanto mais dolorosas se fazem sentir as divisões da Igreja que rompem a participação comum na mesa do Senhor, tanto mais prementes são as orações pela unidade plena de todos os que creem em Cristo.

No que se refere à intercomunhão, a Igreja Católica reconhece que as Igrejas orientais não em plena comunhão com ela ``têm verdadeiros sacramentos e principalmente, em virtude da sucessão apostólica, o sacerdócio e a Eucaristia'', sendo possível, em circunstâncias oportunas e com aprovação da autoridade eclesiástica, uma certa partilha sacramental. Quanto às comunidades eclesiais saídas da Reforma, a intercomunhão eucarística não é possível para a Igreja Católica porque tais comunidades não conservaram ``a genuína e íntegra substância do mistério eucarístico, sobretudo por causa da falta do sacramento da Ordem''. Em caso de grave necessidade, porém, ministros católicos podem administrar os sacramentos da Eucaristia, Penitência e Unção dos Enfermos a cristãos não católicos que os pedem com fé católica e nas devidas disposições.

\section{Penhor da glória futura (1402--1405)}

Numa antiga oração, a Igreja aclama assim o mistério da Eucaristia: ``Ó sagrado banquete, em que se recebe Cristo e se comemora a sua Paixão, em que a alma se enche de graça e nos é dado o penhor da futura glória''. A Eucaristia é, portanto, ao mesmo tempo memória do passado --- a Paixão de Cristo ---, fonte do presente --- a graça que preenche a alma --- e antecipação do futuro --- o penhor da glória celeste. Se a Eucaristia é o memorial da Páscoa do Senhor, se pela Comunhão no altar somos cumulados da ``plenitude das bênçãos e graças do céu'', ela é também a antecipação da glória eterna.

Na última Ceia, o próprio Senhor chamou a atenção dos seus discípulos para a consumação da Páscoa no Reino de Deus: ``Eu vos digo que não voltarei a beber deste fruto da videira, até o dia em que o beberei convosco no Reino do meu Pai''. Sempre que a Igreja celebra a Eucaristia, lembra-se desta promessa e o seu olhar volta-se para ``Aquele que vem''; na sua oração, ela clama pela sua vinda: ``Marana tha --- Vem, Senhor Jesus!'' A Igreja sabe que o Senhor vem já na sua Eucaristia e que está ali no meio de nós, mas que esta presença é velada; por isso, celebra a Eucaristia ``enquanto aguarda a feliz esperança e a vinda de Jesus Cristo nosso Salvador''.

Desta grande esperança --- dos novos céus e da nova terra onde habitará a justiça --- não temos garantia mais segura nem sinal mais manifesto do que a Eucaristia. Com efeito, cada vez que se celebra este mistério, ``realiza-se a obra da nossa redenção'' e nós ``partimos o mesmo pão que é remédio de imortalidade, antídoto para não morrer, mas viver em Jesus Cristo para sempre''. A Eucaristia inscreve, assim, a existência cristã numa tensão escatológica irrenunciável: o cristão que comunga não vive apenas no presente, mas já habita a anterioridade do que ainda não veio, interpelado pela plenitude do Reino que a Eucaristia faz resplandecer em cada celebração.

\section{Em síntese (1406--1419)}

Jesus diz: ``Eu sou o pão vivo descido do céu. Quem comer deste pão viverá eternamente; quem come a minha Carne e bebe o meu Sangue tem a vida eterna, permanece em Mim e Eu nele''. A Eucaristia é o coração e o cume da vida da Igreja porque, nela, Cristo associa a sua Igreja e todos os seus membros ao seu sacrifício de louvor e de ação de graças, oferecido ao Pai de uma vez por todas na Cruz; por este sacrifício, Ele derrama as graças da salvação sobre o seu Corpo, que é a Igreja.

A celebração eucarística inclui sempre: a proclamação da Palavra de Deus; a ação de graças a Deus Pai por todos os seus benefícios, sobretudo pelo dom do seu Filho; a consagração do pão e do vinho; e a participação no banquete litúrgico pela recepção do Corpo e do Sangue do Senhor. Estes elementos constituem um só e mesmo ato de culto. A Eucaristia é o memorial da Páscoa de Cristo, isto é, da obra de salvação realizada pela vida, morte e ressurreição de Cristo, tornada presente pela ação litúrgica. O sacrifício eucarístico é oferecido também em reparação dos pecados dos vivos e dos defuntos e para obter de Deus benefícios espirituais ou temporais.

É o próprio Cristo, Sumo e Eterno Sacerdote da Nova Aliança, quem, agindo pelo ministério dos sacerdotes, oferece o sacrifício eucarístico; e é o mesmo Cristo, realmente presente sob as espécies do pão e do vinho, a oferenda do sacrifício. Só os sacerdotes validamente ordenados podem presidir à Eucaristia e consagrar o pão e o vinho para que se tornem o Corpo e o Sangue do Senhor. Os sinais essenciais do sacramento eucarístico são o pão de trigo e o vinho da videira, sobre os quais é invocada a bênção do Espírito Santo e o sacerdote pronuncia as palavras da consagração ditas por Jesus na última Ceia.

Aquele que quiser receber Cristo na Comunhão eucarística deve encontrar-se em estado de graça; se alguém tiver consciência de pecado mortal, não deve aproximar-se da Eucaristia sem antes ter recebido a absolvição no sacramento da Penitência. A sagrada Comunhão do Corpo e Sangue de Cristo aumenta a união do comungante com o Senhor, perdoa-lhe os pecados veniais e preserva-o dos pecados graves. A Igreja recomenda vivamente aos fiéis que recebam a sagrada Comunhão quando participam na Eucaristia, e impõe-lhes a obrigação de o fazer ao menos uma vez por ano. Tendo passado deste mundo para o Pai, Cristo deixou-nos na Eucaristia o penhor da glória junto d'Ele: a participação no santo sacrifício identifica-nos com o seu Coração, sustenta as nossas forças ao longo da peregrinação desta vida, faz-nos desejar a vida eterna e nos une desde já à Igreja do céu, à Santíssima Virgem e a todos os santos.
